\chapter[Band 5: Funktionen]{Band 5 auf einen Blick}
\noindent\textbf{Funktionen}\par\medskip
\noindent\emph{Lesart.} Sofern eine Zeile nichts anderes angibt,
ist \(F\colon A\to B\) eine Funktion, \(C,D\subseteq A\) und
\(S,T\subseteq B\). Funktionsgraphen werden als Mengen geordneter Paare
verstanden.
\begin{BandUebersicht}
\textbf{Funktion}
\UebFormel{\begin{aligned}
&F\colon A\to B:\quad\TotRel{F,A,B},\\
&(x,y)\in F\dsep(x,z)\in F\vdash y=z.
\end{aligned}}
Zur Totalität tritt die funktionale Eindeutigkeit.
\UebQuellen{\FormulaRefAuto{Funktion}[def];
\FormulaRefAuto{(x,y)\in F\dsep (x,z)\in F \vdash y=z}[ax]}
&
\textit{Eindeutiger Wert und eindeutiger Definitionsbereich}
\UebFormel{\begin{aligned}
&x\in A\vdash\exists!y\;(x,y)\in F,\\
&F\colon A\to B\dsep F\colon C\to D\vdash A=C.
\end{aligned}}
\UebQuellen{\FormulaRefAuto{F\colon A\to B\dsep x\in A\vdash \exists! y\;(x,y)\in F}[thm];
\FormulaRefAuto{FunctionDomainUnique}[thm]}\\
\addlinespace[.8ex]
\textbf{Funktionswert und Graph}
Für \(x\in A\) ist \(F(x)\) das eindeutig bestimmte \(y\) mit
\((x,y)\in F\).
\UebFormel{\Graph(F)\coloneq
\{(x,y)\in A\times B\mid y=F(x)\}.}
\UebQuellen{\FormulaRefAuto{\Graph(F)\coloneq\{(x,y) \in A \times B \mid y=F(x)\}}[def]}
&
\textit{Typisierung und Funktionsextensionalität}
\UebFormel{\begin{aligned}
&x\in A\vdash F(x)\in B,\\
&F,G\colon A\to B\dsep\\
&\quad\forall x\in A\,(F(x)=G(x))\vdash F=G.
\end{aligned}}
\UebQuellen{\FormulaRefAuto{F\colon A\to B\dsep x\in A\vdash F(x)\in B}[thm];
\FormulaRefAuto{FunctionExtensionality}[thm]}\\
\addlinespace[.8ex]
\textbf{Bildmenge}
\(F[C]\) wird durch eindeutige Beschreibung eingeführt:
\UebFormel{y\in F[C]\leftrightarrow\exists x\in C\,(y=F(x)).}
\UebQuellen{\FormulaRefAuto{C\subseteq A\dsep F\colon A\to B\vdash F[C]\coloneqq \iota D\Bigl(\forall y\,\bigl(y\in D \leftrightarrow \exists x\in C\,y=F(x)\bigr)\Bigr)}[def]}
&
\textit{Existenz, Nichtleerheit und Vereinigungen}
\UebFormel{\begin{aligned}
&\exists!Y\,\forall y\,\bigl(y\in Y\\
&\qquad\leftrightarrow\exists x\in C\,(y=F(x))\bigr),\\
&C\in\PowNE{A}\vdash F[C]\in\PowNE{B},\\
&F[C\cup D]=F[C]\cup F[D].
\end{aligned}}
\(\PowNE{A}\) bezeichnet die nichtleeren Teilmengen von \(A\).
\UebQuellen{\FormulaRefAuto{SubsetFunctionImageSetUniqueExists}[thm];
\FormulaRefAuto{NonemptyDirectImage}[thm];
\FormulaRefAuto{F\colon M\to N \dsep A\subseteq M \dsep B\subseteq M \vdash F[A\cup B]=F[A]\cup F[B]}[thm]}\\
\addlinespace[.8ex]
\textbf{Urbild und Faser}
\UebFormel{\begin{aligned}
F^{-1}[T]&\coloneqq\{x\in A\mid F(x)\in T\},\\
\Fib_F(y)&\coloneqq F^{-1}[\{y\}]\quad(y\in B).
\end{aligned}}
\UebQuellen{\FormulaRefAuto{F\colon A\to B\dsep C\subseteq B \vdash F^{-1}[C] \coloneqq \{\,x\in A \mid F(x)\in C\,\}}[def];
\FormulaRefAuto{FiberMapDef}[def]}
&
\textit{Urbildkriterium und Rechengesetze}
\UebFormel{\begin{aligned}
&\forall x\in A\,(F(x)\in T\leftrightarrow x\in C)\\
&\hspace{10mm}\vdash F^{-1}[T]=C,\\
&F^{-1}[S\cup T]=F^{-1}[S]\cup F^{-1}[T],\\
&F^{-1}[S\cap T]=F^{-1}[S]\cap F^{-1}[T].
\end{aligned}}
Die Fasern zerlegen den Definitionsbereich nach gleichen Werten;
leere Fasern sind bei allgemeinen Funktionen erlaubt.
\UebQuellen{\FormulaRefAuto{PreimageEqFromPointwiseIff}[thm];
\FormulaRefAuto{F\colon A\to B \dsep C\subseteq B \dsep D\subseteq B \vdash F^{-1}[C\cup D]=F^{-1}[C]\cup F^{-1}[D]}[thm];
\FormulaRefAuto{F\colon A\to B \dsep C\subseteq B \dsep D\subseteq B \vdash F^{-1}[C\cap D]=F^{-1}[C]\cap F^{-1}[D]}[thm]}\\
\addlinespace[.8ex]
\textbf{Funktion aus einem Graphen}
Eine Relation \(R\subseteq A\times B\) kann eine Funktion beschreiben,
wenn ihre Werte für jedes Argument eindeutig existieren.
&
\textit{Konstruktionsprinzip}
\UebFormel{\begin{aligned}
&R\subseteq A\times B\dsep\\
&\forall x\in A\,\exists!y\in B\,((x,y)\in R)\\
&\hspace{10mm}\vdash R\colon A\to B.
\end{aligned}}
\UebQuellen{\FormulaRefAuto{UniqueValuedGraphFunction}[thm]}\\
\addlinespace[.8ex]
\textbf{Konstante und Fallabbildung}
Für \(b\in B\), \(f\colon C\to B\) und \(g\colon D\to B\):
\UebFormel{\begin{aligned}
&\ConstFun{A}{B}{b}(x)\coloneqq b\quad(x\in A),\\
&\CaseFun{C}{D}{f}{g}(x)\coloneqq\\
&\qquad\IfThenElse{x\in C}{f(x)}{g(x)}.
\end{aligned}}
Die zweite Formel gilt auf \(C\cup D\).
\UebQuellen{\FormulaRefAuto{Konstante Funktion}[def];
\FormulaRefAuto{CaseFunctionDef}[def]}
&
\textit{Elementarer Graph und Auswertung}
\UebFormel{\begin{aligned}
&b\in B\vdash\{(a,b)\}\colon\{a\}\to B,\\
&b\in B\vdash\{(a,b)\}(a)=b,\\
&f\colon A\to C\dsep g\colon B\to C\dsep x\in A\\
&\hspace{6mm}\vdash\CaseFun{A}{B}{f}{g}(x)=f(x).
\end{aligned}}
\UebQuellen{\FormulaRefAuto{SingletonGraphFunction}[thm];
\FormulaRefAuto{SingletonGraphValue}[thm];
\FormulaRefAuto{f\colon A\to C\dsep g\colon B\to C\dsep x\in A \vdash \CaseFun{A}{B}{f}{g}(x)=f(x)}[thm]}\\
\addlinespace[.8ex]
\textbf{Einschränkung}
\UebFormel{\begin{aligned}
&F\restriction_C\colon C\to B,\\
&(F\restriction_C)(x)\coloneqq F(x)\quad(x\in C).
\end{aligned}}
\UebQuellen{\FormulaRefAuto{RestrictionDef}[def]}
&
\textit{Gleichheitskriterium}
Für \(G\colon C\to B\):
\UebFormel{\begin{aligned}
F\restriction_C=G\ \leftrightarrow\ 
\forall x\in C\,(F(x)=G(x)).
\end{aligned}}
\UebQuellen{\FormulaRefAuto{FunctionRestrictionEqualityCriterion}[thm]}\\
\addlinespace[.8ex]
\textbf{Komposition}
Für \(F\colon A\to B\), \(G\colon B\to C\):
\UebFormel{\begin{aligned}
&G\circ F\colon A\to C,\\
&(G\circ F)(x)\coloneqq G(F(x))\quad(x\in A).
\end{aligned}}
\UebQuellen{\FormulaRefAuto{CompositionDef}[def]}
&
\textit{Bilder unter einer Linksinversen}
Für \(F\colon A\to B\), \(G\colon B\to A\), \(K\subseteq A\):
\UebFormel{\begin{aligned}
&\forall x\in A\,G(F(x))=x\\
&\hspace{12mm}\vdash G[F[K]]=K.
\end{aligned}}
\UebQuellen{\FormulaRefAuto{LeftInverseImages}[thm]}\\
\addlinespace[.8ex]
\textbf{Induzierte Potenzmengenabbildung}
\UebFormel{\begin{aligned}
&\PowMap{F}\colon\powerset(A)\to\powerset(B),\\
&\PowMap{F}(X)\coloneqq F[X],\\
&\PowMapNE{F}\coloneqq\PowMap{F}\restriction_{\PowNE{A}}.
\end{aligned}}
\UebQuellen{\FormulaRefAuto{PowerMapDef}[def];
\FormulaRefAuto{NonemptyPowerMapDef}[def]}
&
\textit{Nichtleere Teilmengen bleiben nichtleer}
\UebFormel{\begin{aligned}
&\PowMapNE{F}\colon\PowNE{A}\to\PowNE{B},\\
&(\PowMap{F})^{-1}[\PowNE{B}]=\PowNE{A}.
\end{aligned}}
\UebQuellen{\FormulaRefAuto{NonemptyPowerMapTyped}[thm];
\FormulaRefAuto{NonemptyPowerMapPreimage}[thm]}\\
\addlinespace[.8ex]
\textbf{Spurabbildung}
Für \(F\colon A\to\powerset(B)\):
\UebFormel{\begin{aligned}
&\TraceMap{F}\colon B\to\powerset(A),\\
&\TraceMap{F}(b)\coloneqq\{x\in A\mid b\in F(x)\}.
\end{aligned}}
\UebQuellen{\FormulaRefAuto{TraceMapDef}[def]}
&
\textit{Gleiche Spuren}
Für \(b,c\in B\), \(x\in A\):
\UebFormel{\begin{aligned}
&\TraceMap{F}(b)=\TraceMap{F}(c)\\
&\hspace{8mm}\vdash b\in F(x)\leftrightarrow c\in F(x).
\end{aligned}}
\UebQuellen{\FormulaRefAuto{TraceMapEqualityMembership}[thm]}\\
\addlinespace[.8ex]
\textbf{Erweiterung und überdeckende Familie}
Für \(a\notin A\), \(b\in B\) erweitert \(\ExtMap{F}{a}{b}\)
die Funktion durch den neuen Wert \(a\mapsto b\).
\UebFormel{\ExtMap{F}{a}{b}(x)\coloneqq
\IfThenElse{x=a}{b}{F(x)}.}
\UebQuellen{\FormulaRefAuto{ExtensionMapDef}[def]}
&
\textit{Vereinigung einer Überdeckung}
Für \(D\colon I\to\powerset(W)\):
\UebFormel{\begin{aligned}
&\forall w\in W\,\exists i\in I\,(w\in D(i))\\
&\hspace{12mm}\vdash\bigcup D[I]=W.
\end{aligned}}
\UebQuellen{\FormulaRefAuto{IndexedSubsetFamilyCoverUnion}[thm]}\\
\end{BandUebersicht}
